\section{Introduction}\label{sec:intro}

The Lieb--Thirring inequality bounds the negative eigenvalues of a
Schr\"odinger operator by an integral of its potential. Write
$t_+=\max\{t,0\}$ and $t_-=\max\{-t,0\}$. On the line, let
$0\le W\in L^{\gamma+1/2}(\R)$ and let $H_{-W}$ be the self-adjoint operator
associated with the quadratic form
\[
 h_{-W}[v]=\int_\R|v'(x)|^2\,dx-\int_\R W(x)|v(x)|^2\,dx,
 \qquad v\in H^1(\R).
\]
For $\gamma>1/2$, this form is closed and bounded below, and its negative
spectrum is discrete. We give the needed form and compactness arguments in
Section~\ref{sec:spectral}. Write
\[
 \Tr(H_{-W})_-^\gamma
   =\sum_{\lambda_j(H_{-W})<0}|\lambda_j(H_{-W})|^\gamma,
\]
counting multiplicities; before the inequality is proved this sum is allowed
to be infinite. The optimal constant is
\[
 L_{\gamma,1}=
 \sup_{\substack{0\le W\in L^{\gamma+1/2}(\R)\\W\not\equiv0}}
 \frac{\Tr(H_{-W})_-^\gamma}{\displaystyle\int_\R W^{\gamma+1/2}}.
\]
Two natural comparisons are the semiclassical constant
\begin{equation}\label{eq:semiclassical}
 L_{\gamma,1}^{\cl}
 =\frac1{2\pi}\int_\R(1-\xi^2)_+^\gamma\,d\xi
 =\frac{\Gamma(\gamma+1)}{2\sqrt\pi\,\Gamma(\gamma+3/2)}
\end{equation}
and the optimal constant when only the lowest eigenvalue is retained,
denoted by $L_{\gamma,1}^{(1)}$. The one-dimensional Lieb--Thirring
conjecture predicts the one-bound-state constant for
$1/2<\gamma<3/2$ and the semiclassical constant for $\gamma\ge3/2$.

\begin{theorem}\label{thm:main}
For every $1/2<\gamma<3/2$ and every nonnegative
$W\in L^{\gamma+1/2}(\R)$,
\begin{equation}\label{eq:main}
 \Tr(H_{-W})_-^\gamma
 \le 2\left(\frac{\gamma-1/2}{\gamma+1/2}\right)^{\gamma-1/2}
       L_{\gamma,1}^{\cl}\int_\R W^{\gamma+1/2}.
\end{equation}
The constant is optimal and equals $L_{\gamma,1}^{(1)}$.
More precisely, if $r=(\gamma-1/2)^{-1}$, equality is attained by
$W(x)=(r+1)\sech^2(rx)$.
\end{theorem}

The estimate includes rough, unbounded potentials and infinitely many
negative eigenvalues. Corollary~\ref{spectral:signed} gives the corresponding
inequality for signed potentials on their natural dense weighted form
domain. The theorem concerns scalar potentials on the line; the matrices
used below encode finite orthonormal families. Its sharpness assertion
exhibits an equality potential without classifying all optimizers.

\subsection{History and relation to prior work}

Lieb and Thirring introduced these spectral inequalities in their work on
the kinetic energy of fermions and the stability of matter
\cite{LiebThirring1975,LiebThirring1976}. At exponent one, optimizing over
the potential makes the spectral inequality equivalent to a lower bound
on the kinetic energy of an orthonormal family in terms of its density.
On the line, for a finite orthonormal family $(u_j)\subset H^1(\R)$, this
density is $\rho=\sum_j|u_j|^2$, and the corresponding kinetic bound
controls $\int\rho^3$ by $\sum_j\int|u_j'|^2$.
The one-state variational problem goes back to Keller \cite{Keller1961}.
The conjecture asks whether allowing arbitrarily many bound states can
improve on its constant in the lower exponent range. We verify the
one-state equality potential and its normalization directly.

The scalar upper endpoint predates the general Lieb--Thirring inequalities.
Inverse-scattering trace identities for the Korteweg--de Vries equation
give the sharp value $L_{3/2,1}=3/16$
\cite{ZakharovFaddeev1971,GardnerGreeneKruskalMiura1974};
Lieb and Thirring explicitly credit Gardner, Greene, Kruskal and Miura
in \cite[Section~4(B)]{LiebThirring1976}.
Aizenman and Lieb's monotonicity argument extends the semiclassical
equality to all larger exponents \cite{AizenmanLieb1978}.
At the other endpoint, Weidl first proved finiteness of $L_{1/2,1}$
\cite{Weidl1996}, and Hundertmark, Lieb and Thomas obtained its sharp
value $1/2$ \cite[Theorem~1]{HundertmarkLiebThomas1998}.
Laptev and Weidl established the operator-valued upper-endpoint estimate
and used dimension lifting to obtain the sharp semiclassical constants
for every dimension and $\gamma\ge3/2$
\cite[Theorems~2.1, 2.2 and~3.1]{LaptevWeidl2000}.

Quantitative estimates and orthonormal-function methods narrowed the
remaining gap. Hundertmark, Laptev and Weidl proved
$L_{\gamma,1}\le2L_{\gamma,1}^{\cl}$ for $1/2\le\gamma<3/2$
\cite[Theorem~4.1]{HLW2000}.
At exponent one, Dolbeault, Laptev and Loss extended the method of
Eden and Foias to matrix-valued potentials
\cite{EdenFoias1991,DLL2008}. Their all-dimensional semiclassical ratio
was $\pi/\sqrt3$; subsequent bounds were $1.456$ by Frank, Hundertmark,
Jex and Nam \cite[Theorem~1]{FHJN2021} and the reported $1.44655$
of Corso and Ried \cite[Corollary~1.7]{CorsoRied2025}.
These bounds preceded the sharp exponent-one theorem described below.
In the full one-dimensional strict interval, Levitt's numerical
investigation observed convergence to the one-state profile or splitting
into increasingly separated copies \cite[Section~4.1]{Levitt2014}.
Such computations support the conjectured value but do not bound the
contribution of every orthonormal family. That uniform bound is the
obstacle addressed by our action inequality.

Read and Schulz proved the case $\gamma=1$ in
\cite[Theorem~1]{ReadSchulz2026}, obtaining the sharp constant
$4/(3\sqrt3\pi)$. Their proof extends Benguria and Loss's scalar
cumulative-mass argument \cite[Section~2]{BenguriaLoss2004} to
orthonormal families through a matrix correction to the kinetic energy.
Their operator-valued extension also gives the semiclassical ratio
$2/\sqrt3$ in every dimension at exponent one
\cite[Theorem~2 and Section~2]{ReadSchulz2026}.
Our proof retains the individual eigenvalue scales in a finite-interval
action inequality. Its main analytic step is a rank-one comparison for
an energy-dependent regularized matrix field. A continuation argument
chooses a lower triangular coefficient matrix and a path joining two
prescribed matrix endpoints. Together these yield the action inequality
for every $1/2<\gamma<3/2$. Theorem~\ref{thm:main} therefore resolves
the remaining open-interval cases of the one-dimensional Lieb--Thirring
conjecture positively; the case $\gamma=1$ is already covered by
Read--Schulz.

\subsection{Proof strategy}

Put $\sigma=\gamma-1/2\in(0,1)$. The central result is a variational
inequality on a bounded interval, independent of any potential.
Fix $I=(x_-,x_+)$. Given $N$ orthonormal real functions
$u_i\in H^1_0(I;\R)$ and positive numbers
$k_i$, let $u=(u_i)_{i=1}^N$ and $K=\diag(k_i)$. We prove that the action
\[
 \mathcal E_K(a)=\int_I\bigl(|a'|^2+a^{\transpose}K^2a
                         -|a|^{2+2/\sigma}\bigr)\,dx
\]
has supremum at least
$\sum_i\int_{-k_i}^{k_i}(k_i^2-s^2)^\sigma\,ds$
among functions $a=Cu$ with $C$ lower triangular.
Section~\ref{sec:action} states this precisely in
Theorem~\ref{thm:action} and explains the role of triangularity.

The difficulty is to obtain this lower bound for an arbitrary orthonormal
family while retaining its individual scales $k_i$. The argument replaces
a direct optimization over $C$ by a path with prescribed matrix endpoints.
Section~\ref{sec:field} constructs a regularized map $B\mapsto M_\delta(B)$
from real symmetric $N\times N$ matrices to themselves, and a scalar
primitive whose change from $K$ to $-K$ gives the required
endpoint integral. Section~\ref{sec:rankone} proves the analytic comparison
that controls the nonlinear action term when $M_\delta(B)=aa^{\transpose}$.
Its integrated tangent map is an average of unitary
conjugations; a trace inequality removes this average from the comparison.

Section~\ref{sec:path} chooses $C$ and a path from $K$ to $-K$ by
continuation modulo row signs. Compactness and the exclusion of zero rows
make the continuation possible even when $C$ loses rank.
Section~\ref{sec:action-limit} combines the path with the primitive:
failure of the rank-one constraint contributes a negative square to the
action identity. Compactness allows the penalty parameter to tend to zero
at fixed regularization. Only then is the regularization removed.

Section~\ref{sec:spectral} returns to the potential. Take the $u_i$ to be Dirichlet eigenfunctions with
eigenvalues $-k_i^2$, ordered so that $k_1\ge\cdots\ge k_N$.
Lower triangularity lets the eigenfunction equations bound the quadratic
part of the action by $\int W|a|^2$. A scalar maximization bounds the whole
action by a multiple of $\int W^{1+\sigma}$. Finally, cutoff trial spaces
pass from bounded intervals to the line, and finite partial sums exhaust
all negative eigenvalues. The sharpness calculation uses the explicit
potential in Theorem~\ref{thm:main}.
